\documentclass[pra,aps,preprint,floatfix]{revtex4-1}
\usepackage{graphicx, xcolor}
\usepackage{amsmath,amssymb}
\usepackage{dsfont}

%\usepackage{hyperref}
%\usepackage{color}
% Custom definitions start
\newcommand{\Hq}{\mathcal{H}}
\newcommand{\heff}{H_{\scriptscriptstyle \rm eff}}
\newcommand{\hs}{\mathcal{H}_{\scriptscriptstyle \rm S}}
\newcommand{\hdr}{\mathcal{H}_{\scriptscriptstyle \rm DR}}
\newcommand{\Lhl}{\mathcal{H}_{\scriptscriptstyle \rm L}}
\newcommand{\hl}{\mathcal{H}_{\scriptscriptstyle \rm L}}
\newcommand{\ul}{U_{\scriptscriptstyle \rm L}}
\newcommand{\hsl}{\mathcal{H}_{\scriptscriptstyle \rm SL}}
\newcommand{\Lhsl}{\mathcal{H}_{\scriptscriptstyle \rm SL}}
\newcommand{\hdd}{\mathcal{H}_{\scriptscriptstyle \rm dd}}
\newcommand{\Lhloc}{\mathcal{H}_{\scriptscriptstyle \rm Loc}}
\newcommand{\hloc}{H_{\scriptscriptstyle \rm loc}}
\newcommand{\nn}{\nonumber}
\newcommand{\tq}{\mathcal{T}}
\newcommand{\fo}{\mathcal{F}}
\newcommand{\wo}{\omega_{\circ}}
\newcommand{\Ur}{U^{\dagger}_{L}}
\newcommand{\tr}{\text{Tr}}
\newcommand{\trl}{\text{Tr}_{ \scriptscriptstyle \rm L}}
\newcommand{\rh}{\rho_{\scriptscriptstyle  S}}
\newcommand{\rl}{\rho_{\scriptscriptstyle \rm L}^{\scriptscriptstyle \rm eq}}
%\newcommand{\sc}{ \scriptscriptstyle \rm sec}
\newcommand{\so}{\mathcal{S}}
\newcommand{\etal}{\textit{et al.\ }}
\newcommand{\sts}{\vert_{t\rightarrow\infty}}
\newcommand{\stoo}{\vert_{t\rightarrow 0}}
\newcommand{\ic}{\vert_{t=0}}
\newcommand{\ket}[1]{\vert #1 \rangle}
\newcommand{\bra}[1]{\langle #1 \vert}
\newcommand{\comment}[1]{ }
\begin{document}
\title{\emph{Supplementary article}: \\
%Discrete time-crystalline phase in periodically driven dissipative dipolar systems}
Emergence and stability of discrete time-crystalline phases in open quantum systems}
\author{Saptarshi Saha}
\email{ss17rs021@iiserkol.ac.in}
\author{Rangeet Bhattacharyya}
\email{rangeet@iiserkol.ac.in}
\affiliation{Department of Physical Sciences, Indian Institute of Science Education and Research Kolkata,\\
Mohanpur -- 741246, West Bengal, India}



\maketitle
\section{Dynamical Equation } 
In this section, we discuss the detailed dynamics of the two-spin dipolar coupled system in the presence of
a thermal bath. Following the recent experiments, the system is subjected to a two-pulse driving protocol
which consists of two non-commuting Hamiltonians acting in two successive periods
\cite{beatrez_critical_2023, choi_observation_2017}. The system-bath coupling strength is relatively weaker
than the periodic drive and dipolar interaction. Hence, we neglect the effect of system-bath coupling in our
dynamics. A schematic diagram is provided for a better understanding of the dynamics of this system.
\begin{figure}[htb]
\includegraphics[width=0.48\linewidth]{Figure-1.pdf} 
\caption{The schematic diagram shows the experiment demonstrating the DTC phase in a dissipative dipolar 
system. The initial state is prepared by using a  $\pi/2$ pulse in the `$y$' direction. Next, the two-pulse protocol begins; it is shown enclosed in the square brackets. The spin-locking pulse is applied for $\tau_1$ 
time, followed by a $\pi + \delta $ rotation pulse applied along `$y$' direction for a duration $\tau_2$. 
The whole sequence, $\tau = \tau_1 + \tau_2$ is repeated for $n$ times before measuring $S(M_x(t))$.}
\label{fig-1}
\end{figure}
\subsection{Dynamical equation during spin-locking pulse} \label{sub1}
The total Hamiltonian of the system is written as,
\begin{eqnarray}\label{ham}
\mathcal{H}(t)&=& \hs^{\circ} + \hl^{\circ}  + \hsl + \hdd + \mathcal{H}_x(t)+\hl(t).
\end{eqnarray} 
The first two terms are the free Hamiltonians of the system.
$\hsl$ represents the weak system-bath coupling. $\hdd$ is the dipolar interaction between the two spins. We 
select only the secular part of the Hamiltonian, assuming that the effect of the nonsecular part is much
smaller. For details, please refer to our earlier work on prethermalization \cite{saha_cascaded_2023}. The analytical 
form of the secular dipolar Hamiltonian is given by
\begin{eqnarray}
\hdd^{\scriptscriptstyle \rm sec} = \omega_{d_0}\left(2 I_z^1 I_z^2 - I_x^1 I_x^2 - I_y^1 I_y^2 \right)
\end{eqnarray}
Here, $\omega_{d_0}$ is the dipolar coupling strength and $I_i = \sigma_i/2, \, \{i \in x,y,z\}$, where $\sigma_i$ is the Pauli spin matrix. The form of the on-resonance periodic drive $\mathcal{H}_x$ is given by,
$\mathcal{H}_x =\sum\limits_{i=1}^2 \omega_1 I_{x}^i  \cos\omega t $. $\hl(t)$ represents the equilibrium thermal 
fluctuations present in the thermal bath. We neglect the contribution from $\hsl$ in our analysis. However,
the fluctuations contribute to drive-induced dissipation and dipolar relaxation.
$\hl(t) = \sum_i f_i(t)\vert \phi_i
\rangle \langle \phi_i \vert$. Here, $\{\vert\phi_i\rangle\}$ are the eigenbases of the environment, and  $f_i$s are the Gaussian stochastic variables with zero mean and standard deviation $\kappa \,
\left(\kappa^2=\frac{1}{\tau_c}\right)$, and $\tau_c$ is defined as the bath correlation time. Fluctuation-regulated 
quantum master equation (FRQME) can be used for such cases to derive the dynamical equation of the system 
\cite{chakrabarti2018b}. In the interaction picture of $ \hs^{\circ} + \hl^{\circ} $, the analytical form of FRQME 
for this case is written as, 
\begin{eqnarray}\label{dynamics-1}
\frac{d\rh}{dt}&=& -i\left[H^{\scriptscriptstyle \rm sec}, \rh(t)\right] -  \tau_c\left[ H^{\scriptscriptstyle \rm sec},\left[H^{\scriptscriptstyle \rm sec},\rh(t) \right] \right],\qquad [0<t<\tau_1]
\end{eqnarray}
Here, $H^{\scriptscriptstyle \rm sec}$ is the interaction representation of the secular part of 
$\mathcal{H}_x(t) + \hdd$. The form of the secular part is given by, $H_x = \sum\limits_{i=1}^2 \omega_1 I_{x}^i$ 
and $H_{\scriptscriptstyle \rm dd} = \omega_{d_0} \left(2 I_z^1 I_z^2 - I_x^1 I_x^2 - I_y^1 I_y^2 \right) $. The 
representation of the symmetric observables is written as
\begin{eqnarray}
M_{\alpha} &=&   {\rm Tr}_s  [ (I_{\alpha}\otimes \mathbb{I} + \mathbb{I} \otimes I_{\alpha})\rho_s ] \nn\\
M_{\alpha\beta}&=&   {\rm Tr}_s [(I_{\alpha} \otimes I_{\beta} + I_{\beta} \otimes I_{\alpha})
\rho_s ],\quad \forall\ \alpha\neq\beta \nn \\
M_{\alpha\alpha} &=&  {\rm Tr}_s [(I_{\alpha}\otimes I_{\alpha})\rho_s]. 
\label{observables}
\end{eqnarray} 
Here, $\alpha, \beta \in \{x,y,z\}$. In terms of those observables, we analyze the above Eq. (\ref{dynamics-1}). Such 
dynamical equations were previously reported by us \cite{saha_cascaded_2023}. In the presence of 
$H^{\scriptscriptstyle \rm sec}$, the dynamical equations can be divided into two subgroups. Each group
contains four observables, $\{M_x, M_{yy}, M_{zz}, M_{yz}\}$ and $\{M_y, M_{z}, M_{xy}, M_{xz}\}$. The
dynamical equations for the first group are given by,
\begin{eqnarray}\label{eq-pre} 
\dot{M}_{x} &=& -\frac{9}{4} \omega_{d_0}^2 \tau_c  M_{x} + 6 \omega_1 \omega_{d_0} \tau_c M_{zz} -6 \omega_1 \omega_{d_0} \tau_c M_{yy} -3\omega_{d_0} M_{yz} \\
\dot{M}_{zz} &=&  \frac{3}{4}\omega_1 \omega_{d_0} \tau_c M_{x} -2\omega_1^2 \tau_c M_{zz} +  2\omega_1^2 \tau_c M_{yy} + \omega_1 M_{yz}  \\
\dot{M}_{yy} &=&  -\frac{3}{4}\omega_1 \omega_{d_0} \tau_c M_{x} +2\omega_1^2 \tau_c M_{zz} -  2\omega_1^2 \tau_c M_{yy} - \omega_1 M_{yz}  \\
\dot{M}_{yz} &=& \frac{3}{4}\omega_{d_0} M_{x} -2 \omega_1 M_{zz} +2\omega_1 M_{yy} -(4 \omega_1^2 + \frac{9}{4}\omega_{d_0}^2 ) \tau_c M_{yz}
\end{eqnarray}
There exist several conserved quantities in this block, which are given by,
\begin{eqnarray}
3\omega_{d_0}\dot{M}_{zz}+\omega_1  \dot{M}_{x} &=& 0\\
\dot{M_{yy}}+\dot{M_{zz}} &=& 0\\
\dot{M}_{xx} &=& 0.
\end{eqnarray}
The existence of such conserved quantities ensures that, if initially, any one of the observables from the first group 
is non-zero, then they have finite values at the steady state. The spin-lock sequence leads to the prethermalization 
in the system \cite{beatrez_floquet_2021,saha_cascaded_2023}. Although, by adding $\hsl$ and other non-secular
contributions, such terms will eventually 
vanish. We are only interested in the dynamics of $M_x$ for a short time compared to the time required for
thermalization (a condition amply satisfied in experiments). The solution 
of $M_x(t)$ from the above Eq. (\ref{eq-pre}) for $M_x \stoo = M_\circ$ is given as,
\begin{eqnarray}\label{mxpre}
M_x^{\scriptscriptstyle \rm pre} (M_\circ, t) = M_\circ \left(\frac{4 \omega_1^2}{\kappa_1^2} + \frac{9}{4}\frac{\omega_{d_0}^2}{\kappa_1^2} \cos (\kappa_1 t) e^{- \kappa_1^2 t \tau_c}\right) 
\end{eqnarray}  
Here, $\kappa_1^2 = 4 \omega_1^2 + \frac{9}{4}\omega_{d_0}^2$.
The above solution (Eq. (\ref{mxpre})) signifies that $M_x(t)$ will reach a steady state, and the steady state value 
is given as $M_x^{\scriptscriptstyle \rm pre}\sts = M_\circ \frac{4 \omega_1^2}{\kappa_1^2}$. This steady
state is a prethermal state \cite{saha_cascaded_2023}. The presence of the conserved quantity also signifies that the 
initial magnetization $(M_\circ)$ is shared between $M_x$ and $M_{zz}$. 
\begin{figure*}[htb]
 %\hspace*{-1mm}
 \raisebox{3cm}{\normalsize{\textbf{(a)}}}\hspace*{-1mm}
\includegraphics[width=0.36\linewidth]{Figure-2s.pdf} 
\hspace*{3mm}
\raisebox{3cm}{\normalsize{\textbf{(b)}}}\hspace*{-1mm}
\includegraphics[width=0.36\linewidth]{Figure-3s.pdf} 
\caption{ Fig. (a) shows the plot of observables $M_x^{\scriptscriptstyle \rm pre}, \,
M_{zz}^{\scriptscriptstyle \rm pre}, \, M_{yy}^{\scriptscriptstyle \rm pre}, M_{yz}^{\scriptscriptstyle \rm
pre}$ versus no. of cycle by numerically solving Eq. (\ref{dynamics-1}). Here $\omega_2 \tau_2 = 0$. The
list of fixed parameters are given as, $\omega_1 = 2\pi \times 40$ KHz, $\omega_{d_0} = 2 \pi \times 10$
KHz, $\tau_c = 10^{-4}$ ms and $\tau_1 = 2\pi \times 0.02 / \omega_1$ KHz$^{-1}$. The initial condition is
chosen as $M_x \stoo =1$. For such a choice of parameters, the system reaches a steady state after some
cycle. The dynamical evolution of $M_{zz}^{\scriptscriptstyle \rm pre}, \, M_{yy}^{\scriptscriptstyle \rm
pre}, M_{yz}^{\scriptscriptstyle \rm pre}$ are negligible compared to $M_x^{\scriptscriptstyle \rm pre}$.
Fig. (b) shows the plot of observables $M_z^{\scriptscriptstyle \rm pre}$ versus the number of cycles for three
different choices of $\tau_1, \tau_c$. Here $\omega_2 \tau_2 = 0$. The list of fixed parameters are
$\omega_1 = 2\pi \times 40$ KHz, $\omega_{d_0} = 2 \pi \times 4$ KHz. Three choices are given as $\{\tau^1_1
= 2\pi \times 0.05 / \omega_1$ KHz$^{-1}$, $\tau^1_c = 10^{-4}$ ms $\}$,  $\{\tau^2_1 = 2\pi \times 0.2 /
\omega_1$ KHz$^{-1}$, $\tau^1_c = 10^{-4}$ ms $\}$, $\{\tau^1_1 = 2\pi \times 0.05 / \omega_1$ KHz$^{-1}$, 
$\tau^2_c = 10^{-3}$ ms $\}$. The initial condition is chosen as $M_z\stoo = -0.2$. The above plot 
indicates that for higher values of $\tau_c$ and $\tau_1$, $M_z^{\scriptscriptstyle \rm pre}$ decays faster 
for increasing the number of cycles. Such a condition is necessary for obtaining the DTC phase.}
\label{fig-2}
\end{figure*}
For the demonstration of the DTC phase, the loss of `$x$' magnetization must be very small in each cycle of the 
evolution. Otherwise, $M_x(t)$ will vanish after a few cycles, so we must have $M_x\stoo /M_x\sts \to 1$. Such condition will 
be satisfied if $\omega_1 > \omega_{d_0} $. Choi \etal have chosen the drive and dipolar coupling in such a way that the 
condition is satisfied for their experiment \cite{choi_observation_2017}. From the above plot (Fig. \ref{fig-2}), the 
evolution of $\{M_{zz}, M_{yy}, M_{yz}\}$ are negligible compared to $M_x$ in the limit $\omega_1 > \omega_d$, so for 
the remaining part of the calculation, we ignore their effect on $M_x(t)$. The dynamical equations for the other group 
are written as,
\begin{eqnarray}\label{eq-pre2} 
\dot{M}_{z} &=& - \omega_{1}^2 \tau_c  M_{z} + \omega_1 M_y + 3 \omega_1 \omega_{d_0} \tau_c M_{xz}  \\
\dot{M}_{y} &=& -\omega_1 M_z -(\omega_1^2 + \frac{9}{4}\omega_{d_0}^2 ) \tau_cM_{y} +3 \omega_{d_0} M_{xz} +  3 \omega_1 \omega_{d_0} \tau_c M_{xy}\  \\
\dot{M}_{xz} &=& \frac{3}{4} \omega_1 \omega_{d_0} \tau_c M_{z} -\frac{3}{4} \omega_{d_0} M_{y}-(\omega_1^2 + \frac{9}{4}\omega_{d_0}^2 ) \tau_c M_{xz} + \omega_1 M_{xy}\\
\dot{M}_{xy} &=& \frac{3}{4} \omega_1 \omega_{d_0} \tau_c M_{y} -\omega_1M_{yz} -\omega_1^2 \tau_c M_{xy}
\end{eqnarray}
As there exist no conserved quantities involving the observables from this group, all the observables will vanish at 
the steady state even if they have any initial non-zero values. There is no closed analytical
form for the observables. Hence we provide the plot of the observable $M_z^{\scriptscriptstyle \rm pre}
(M_\alpha, t)$ by numerically solving Eq. (\ref{eq-pre2})  for a different choice of $\tau_1, \, \tau_c$.
Here, $M_z\stoo = M_\alpha$. Other observables are not relevant for this case.
\subsection{Dynamical equation for rotation along the `$y$' direction}\label{sub2}
Next, we consider the rotation along the `$y$' direction. As this is a short-duration pulse, no heating occurs in this 
time duration. In other words, the effects of the dissipators are negligible during this period. Therefore, the dynamics 
is adequately described by the first-order process. The dynamical equation in the interaction frame is given by
\begin{eqnarray}\label{dynamics-2}
\frac{d\rh}{dt}&=& -i\left[H_y, \rh(t)\right],
\end{eqnarray}
Here, $H_y = \omega_2 \sum_{i=1}^2 I_y^i$. 
To obtain the existence of the DTC phase, we are only interested in the dynamics of $M_x$. The solution of Eq. (\ref{dynamics-2}) in terms of $M_x$ is written as,
\begin{eqnarray}
M_x^{\scriptscriptstyle \rm rot} (t) = M_\circ \cos(\omega_2 t),\quad M_z^{\scriptscriptstyle \rm rot} (t) = M_\circ \sin(\omega_2 t) 
\end{eqnarray}
Here $M_x^{\scriptscriptstyle \rm rot} \stoo = M_\circ$. 
\section{Dynamical evolution upto $2 \tau$ time-period}
Using the dynamical equation in the subsections \ref{sub1} and \ref{sub2},  we provide the analytical solution of the 
time evolution of $M_x(t)$ in this section. The dynamical equation in the Liouville space is written as,
\begin{eqnarray}\label{main-eq}
\hat{\rh}(t) =\left[e^{\hat{\mathcal{L}}_y  \tau_2} e^{\hat{\mathcal{L}}_{\scriptscriptstyle \rm sec}  \tau_1}\right]^n \hat{\rh} \stoo.
\end{eqnarray}
Here, $\hat{\mathcal{L}}_y$ is the Liouvillian corresponding to $H_y$. We also define, $\hat{\mathcal{L}}_1$ is the 
Liouvillian corresponding to $H_{\scriptscriptstyle \rm dd} + H_x$, therefore, $\hat{\mathcal{L}}_{\scriptscriptstyle \rm sec} 
= \hat{\mathcal{L}}_1 + \tau_c \hat{\mathcal{L}}_1 \times\hat{\mathcal{L}}_1 $. The time period $(\tau)$ is defined as, 
$\tau = \tau_1 + \tau_2$. For the initial condition $M_x \stoo = 1$, the solution of $M_x(t)$ for 
$\omega_1> \omega_{d_0}$ at time $\tau_1$ is given as, 
\begin{eqnarray}
M_x (\tau_1) = M_x^{\scriptscriptstyle \rm pre}(1, \tau_1),\quad M_z (\tau_1) = 0.
\end{eqnarray}
After the rotation, the solution is given by,
\begin{eqnarray}
M_x (\tau) = M_x^{\scriptscriptstyle \rm pre}(1, \tau_1) \cos\theta, \quad  M_z (\tau) = M_x^{\scriptscriptstyle \rm pre}(1, \tau_1) \sin\theta
\end{eqnarray}
Here, $\theta = \omega_2 \tau_2$. Similarly, after a time $ \tau + \tau_1$, the solution is,
\begin{eqnarray}
M_x (\tau + \tau_1) = M_x^{\scriptscriptstyle \rm pre}(M_x^{\scriptscriptstyle \rm pre}(1, \tau_1), \tau_1) \cos\theta , \quad M_z (\tau + \tau_1) = M_z^{\scriptscriptstyle \rm pre}(M_x^{\scriptscriptstyle \rm pre}(1, \tau_1), \tau_1) \sin \theta
\end{eqnarray}
Finally after the $2\tau$ time-period, the analytical expression of $M_x(2\tau)$ is written as,
\begin{eqnarray}
M_x(2 \tau) = M_x^{\scriptscriptstyle \rm pre}(M_x^{\scriptscriptstyle \rm pre}(1, \tau_1), \tau_1) \cos^2 \theta + M_z^{\scriptscriptstyle \rm pre}(M_x^{\scriptscriptstyle \rm pre}(1, \tau_1), \tau_1) \sin^2 \theta 
\end{eqnarray}
\subsection{Solution for $\theta = \pi$}
Putting $\theta = \pi$, in the above solutions, we get,
\begin{eqnarray}\label{eqpi}
M_x (0) &=&1,\\
M_x (\tau_1) &=& M_x^{\scriptscriptstyle \rm pre}(1, \tau_1),\\
M_x (\tau) &=& -M_x^{\scriptscriptstyle \rm pre}(1, \tau_1),\\
M_x (\tau + \tau_1) &=& -M_x^{\scriptscriptstyle \rm pre}(M_x^{\scriptscriptstyle \rm pre}(1, \tau_1), \tau_1)\\
M_x (2\tau) &=& M_x^{\scriptscriptstyle \rm pre}(M_x^{\scriptscriptstyle \rm pre}(1, \tau_1), \tau_1) 
\end{eqnarray}
In the above Eq. (\ref{eqpi}), we show that, in every time-period $\tau$, the sign of $M_x$ is reversed.
Therefore, $M_x(t)$ comes to the same phase at $2\tau$ time-period. Therefore the system shows the 
period-doubling response, which exhibits the notion of the DTC phase in the system. On the other hand, 
$M_z(t)$ remains zero throughout the evolution.
\begin{figure*}[htb]
\includegraphics[scale =0.9]{Figure-4s.pdf} 
\caption{ Plot of $M_x$ versus the number of cycles is shown in  (a), (c),  (e) and their corresponding Fourier 
transforms, $S(M_x)$ versus $\omega$ are shown in  (b), (d), and (f). The chosen parameters are:
$\omega_1 = 2\pi \times 50$ KHz, $\omega_2 = 2\pi \times 100$ KHz, $\omega_{d_0} = 2\pi \times 0.1$ KHz, 
$n=200,$ $\tau_1 = 2\pi\times 0.02/ \omega_1$ KHz$^{-1}$ and $\tau_2 = 1.04\pi/ \omega_2$ KHz$^{-1}$. 
The initial state is chosen as, $\vert\psi\rangle \stoo = \prod_{i=1}^2 \left(\vert\uparrow \rangle_i 
+ \vert\downarrow \rangle_i\right)/ \sqrt{2}$. For the upper plot (Fig. 2 (a), (b); red color online), 
the value of $\tau_c = 10^{-3}$ ms. For the middle plot (Fig. 2 (c), (d); blue color online), the 
value of $\tau_c = 10^{-5}$ ms.  For the lower plot (Fig. 2 (e),(f); green color online), the value 
of $\tau_c = 10^{-7}$ ms. The above plots signifies that lowering $\tau_c$ results in the decaying of 
robust $2 \tau$ response in the system.}
\label{fig-3}
\end{figure*}

\subsection{Solutions for $\theta = \pi + \delta$, with $\delta/ \pi \to 0$}
In such cases, we can assume that, $\cos(\pi + \delta ) = -1$, $\sin(\pi + \delta ) = -\delta$. The solution is given by,
\begin{eqnarray}\label{eqdelta}
M_x (0) &=&1,\\
M_x (\tau_1) &=& M_x^{\scriptscriptstyle \rm pre}(1, \tau_1), \quad M_z (\tau_1) = 0,\\
M_x (\tau) &=& -M_x^{\scriptscriptstyle \rm pre}(1, \tau_1), \quad  M_z (\tau) = -\delta M_x^{\scriptscriptstyle \rm pre}(1, \tau_1)\\
M_x (\tau + \tau_1) &=& -M_x^{\scriptscriptstyle \rm pre}(M_x^{\scriptscriptstyle \rm pre}(1, \tau_1), \tau_1), \quad M_z (\tau + \tau_1) = -\delta M_z^{\scriptscriptstyle \rm pre}(M_x^{\scriptscriptstyle \rm pre}(1, \tau_1), \tau_1)\\
M_x(2 \tau) &=&  M_x^{\scriptscriptstyle \rm pre}(M_x^{\scriptscriptstyle \rm pre}(1, \tau_1), \tau_1)  + \delta^2 M_z^{\scriptscriptstyle \rm pre}(M_x^{\scriptscriptstyle \rm pre}(1, \tau_1), \tau_1)
\end{eqnarray}
In the above solutions (Eq. (\ref{eqdelta})), the $2\tau$ periodicity is absent due to the presence of the 
extra term in $M_x (2 \tau)$. Such a response can be retrieved if $\delta$ is small. From Fig. \ref{fig-2}(b), 
it is clear that, $M_z(M_\alpha,t)$ decays faster for higher values of $\tau_1$. Therefore, higher $\tau_1$ and 
lower $\delta$ are suitable conditions for observing the DTC phase in the system.
\subsection{$\tau_c$ dependency of the DTC phase}
In the limits of $\omega_1 > \omega_d$, the expressions of $M_z^{\scriptscriptstyle \rm pre} (M_\alpha,t)$ is given as,
\begin{eqnarray}
M_z^{\scriptscriptstyle \rm pre} (M_\alpha, t) =  M_\alpha \cos (\omega_1 t)  e^{-\omega_1^2  \tau_c t}
\end{eqnarray}
So, for higher values of $\tau_c$, the effect of the extra term of $M_x (2\tau)$ in Eq. (\ref{eqdelta}) can be 
neglected so, $2\tau$ periodicity is robust in this regime. In the Fig. \ref{fig-2}(b), we show that 
$M_z^{\scriptscriptstyle \rm pre} (M_\alpha, t)$ decays faster for higher values of $\tau_c$, which is
helpful for retrieving the DTC regime. We also plot $M_x (t)$ and it's Fourier transform $\mathcal{S}(M_x (t))$. 
For low values of $\tau_c$, there is no existence of $2\tau$ periodicity, but for changing the value of $\tau_c$ 
from $10^{-7}$ ms to $10^{-3}$ ms, such a novel response is retrieved.


\bibliographystyle{apsrev4-1}
\bibliography{references}

\end{document}








